

%قبل از شروع مسئله ترجیحاً چیزی ننویسید. اما اگر خواستید بنویسید، در قالب لیستی (فقرات و شمارش) بنویسید.
%برای اضافه کردن نمره، بعدی از شروع‌مسئله، [x] را اضافه کنید و به جای x  نمره مورد نظر را قرار دهید. برای زیرمسئله‌ها هم به همین صورت است.
%\شروع‌مسئله[10]
\شروع‌مسئله []
%\متن‌سیاه{سوال عملی اول:}\\
%\متن‌سیاه{دراپ‌تست}\\

شنگدباو $n$ تا قفس دارد که در هر کدام موشی قرار دارد. تعدادی جفت از این قفس‌ها به صورت یکطرفه به هم متصل اند و موشی که در قفس اول باشد می‌تواند به قفس دوم در آن جفت برود. 

رفتار موش‌های شنگدباو عجیب است! موش‌های او اگر در قفسی باشند که به $i$ تا قفس دیگر راه‌داشته باشد به احتمال 
$\frac{1}{i}$
به یکی از قفس‌های دیگر می‌روند. فرض کنید شنگدباو موش‌ها را در این قفس‌ها با راه‌های بینشان رها کرده‌است و $N$ روز بعد می‌خواهد برگردد و موش‌ها را جمع‌آوری کند. می‌دانیم $N$ عدد بزرگی است و می ‌توانید فرض کنید
$N \to \infty$.

سؤال اینجاست که پس از این تعداد روز شنگدباو حداقل چند‌تا قفس را باید بردارد تا تضمین شود به احتمال حدوداً «۱» همه موش‌ها را جمع آوری کرده‌است.

راهنمایی: از مؤلفه‌های قویاً همبند گراف قفس‌ها و راه‌ها استفاده کنید و ببینید پس از تعداد زیادی روز موش‌ها در کدام مؤلفه‌ها ممکن است جمع شوند.

\متن‌سیاه{ورودی:}\\
در خط اول ورودی به‌شما $n$ و $m$ داده می‌شود که بیانگر تعداد قفس‌ها و تعداد راه‌های یکطرفه بینشان است. سپس در $m$ خط بعدی دو عدد $u, v$ داده‌می‌شود که به معنای وجود مسیری از قفس $u$ به قفس $v$ است.

$$1 \le n, m \le 10^5$$

تضمین می‌شود راهی از یک قفس به خودش نیست. 

تضمین می‌شود دو تا جفت برابر در $m$ خط ورودی داده‌ نمی‌شود. بعبارتی بین هر دوتایی 
$(u, v)$
حداکثر یک‌راه وجود دارد.

\متن‌سیاه{خروجی:}\\
تنها یک عدد خروجی دهید و آن‌هم برابر تعداد کمینه‌ قفس‌هایی باشد که با جمع کردنشان شنگدباو همه موش‌ها را جمع‌آوری می‌کند.

\متن‌سیاه{ورودی و خروجی نمونه:}\\

\begin{latin}
	\begin{tabularx}{\textwidth}{|X | X|}
		\hline
		\textbf{stdin} &
		\textbf{stdout}\\
		\hline
		5 5 & 2\\
		1 2 &\\
		1 3& \\
		2 4&\\
		3 4&\\
		2 5&\\
		\hline
	\end{tabularx}
\end{latin}

\solution{
اگر گراف فاقد دور باشد دارای تعدادی رأس 
\lr{sink}
یعنی بدون خروجی است. به‌سادگی می‌توان دید جواب مسئله برابر تعداد این رئوس است یعنی پس از $n$ مرجله موش‌ها نهایتاً در یکی از این رئوس می‌روند. برای اثبات چنین چیزی ترتیب توپولوژیک گراف را در نظر بگیرید می‌توان دید پس از مرحله $i$ ام اگر رأس $i$ام یک 
\lr{sink}
نباشد خالی می‌شود و دیگر پر نمی‌شود. در نهایت کافیست روی رئوس پیمایش کنیم و تعداد رئوس بدون خروجی را بشماریم.

اما به ازای حالت کلی باید مؤلفه‌های قویا همبند را محاسبه کرد و به جای تک رأس یکسری مؤلفه‌ها در 
\lr{sink}
قرار دارند. می‌توان اثبات کرد پس از بی‌نهایت مرحله احتمال حضور موش در مؤلفه‌های غیر
\lr{sink}
برابر صفر است. شهود آن این است که هر سری به احتمال کمی ممکن است موشی از این مؤلفه خارج شود و به مؤلفه بعدی در این ترتیب توپولوژیک برود، ممکن است این احتمال کم باشد اما به ازای بی‌نهایت مرحله می‌توان دید احتمال خضور موش‌ها در اینجا به صفر میل می‌کند.(برای اثبات دقیق می‌توانید در مورد زنجیر‌های مارکف مطالعه کنید!).

حالا برای حل مسئله باید با استفاده از الگوریتم پیدا کردن مؤلفه‌های قویاً همبند مؤلفه‌های گراف را پیدا کنید و ببینید چندتا مؤلفه 
\lr{sink}
دارد، هر رأسی داخل این مؤلفه‌ها احتمال حضور موش را دارد!

}
%}
% صورت اصلی مسئله را در این‌جا بنویسید.
%اگر مسئله چند قسمت داشت به این صورت عمل کنید:
%\شروع{زیرمسائل}
%\زیرمسئله[0] % مسئلۀ ۱آ
%\زیرمسئله % مسئلۀ ۱ب
%\زیرمسئله % مسئلۀ ۱پ
%\زیرمسئله % مسئلۀ ۱ت
%\پایان{زیرمسائل}



%%%%%%%%%%%%%%%%%%%%%%%%%%%%%%%%%%%%%%%%%%%%%%%%%%%%%
% برخی دستورات مفید برای تمرین‌ها. %
%%%%%%%%%%%%%%%%%%%%%%%%%%%%%%%%%%%%%%%%%%%%%%%%%%%%%

% فرمول‌های نمونه:
%$f(x) = \Theta(x)$
%$T(x, y) \leq \log(x) + 2^y + \binom{2n}{n}$
% {\tt code\_function}


% لیست موارد بدون شماره‌گذاری:
%\شروع{فقرات}
%    \فقره مورد اول 
%        \شروع{فقرات}
%            \فقره مورد اول داخلی 
%                \شروع{فقرات}
%                    \فقره اولین مورد خیلی داخلی 
%                    \فقره دومین مورد خیلی داخلی 
%                \پایان{فقرات}
%            \فقره مورد دوم داخلی 
%        \پایان{فقرات}
%    \فقره مورد دوم 
%\پایان{فقرات}

% لیست موارد شماره‌گذاری شده:
%\شروع{شمارش}
%    \فقره یک مورد 
%    \فقره یک مورد دیگر 
%    \فقره باز هم یک مورد دیگر
%\پایان{شمارش}

% برای نوشتن یک فرمول چند خطی که یک شماره بخوره. برای حالت‌های دیگر به سایت زیر مراجعه کنید:
% https://www.overleaf.com/learn/latex/Aligning_equations_with_amsmath
%\begin{align} 
%    \begin{split}
%        (x+y)^3 &= (x+y)^2(x+y) \\
%                &= (x^2+2xy+y^2)(x+y) \\
%                &= (x^3+2x^2y+xy^2) + (x^2y+2xy^2+y^3) \\
%                &= x^3+3x^2y+3xy^2+y^3
%    \end{split}                                 
%\end{align}

% نوشتن ماتریس و آرایه دوبعدی
%\begin{align}
%    A = 
%    \begin{bmatrix}
%        A_{11} & A_{21} \\
%        A_{21} & A_{22}
%    \end{bmatrix}
%\end{align}

% اضافه کردن تصویر با شرح. برای اضافه کردن تصویر بدون شرح، فقط دستور درج‌تصویر را استفاده کنید.
%\شروع{شکل}[h!]
%\تنظیم‌ازوسط
%  \درج‌تصویر[پهنا=0.5\پهنای‌متن]{img.jpg}
%\شرح{شرح تصویر}
%\برچسب{تصویر:۱}
%\پایان{شکل}
%\پایان{زیرمسائل}

%کد زیر را ببینید.
%%\lstinputlisting{source_filename.py}
%\begin{latin}
%\begin{lstlisting}
%# Import the modules
%import sys
%import random
%
%ans = True
%
%while ans:
%question = raw_input("Ask the magic 8 ball a question: (press enter to quit) ")
%
%answers = random.randint(1,8)
%
%if question == "":
%sys.exit()
%
%elif answers == 1:
%print "It is certain"
%
%elif answers == 2:
%print "Outlook good"
%
%elif answers == 3:
%print "You may rely on it"
%
%elif answers == 4:
%print "Ask again later"
%
%elif answers == 5:
%print "Concentrate and ask again"
%
%elif answers == 6:
%print "Reply hazy, try again"
%
%elif answers == 7:
%print "My reply is no"
%
%elif answers == 8:
%print "My sources say no"
%\end{lstlisting}
%\end{latin}

	
	